% s-184-p184-m001; original PDF 184
\Aphorism{s-184-p184-m001}{273}
Mancher sagt die Wahrheit, sagt sie aber mit so vielen Tact-, Accent-, Tempofehlern, dass er gegen sie und sich einnimmt. Auch die Logik gestattet Grazie, auch mit Thatbeständen lässt sich Musik machen; ja gerade für das Unwiderleglichste, weil es herrisch auftritt und verletzt, muss Liebenswürdigkeit des Vortrags die wärmste Fürsprache einlegen. Ein Gedanke, der befehlen will — und das will ja die „Wahrheit“ — muss sich jeden Augenblick durch vollendete Höflichkeit legitimiren. Gedanken, die nur „zu Besuch kommen“ — gemeinhin Hypothesen, Einfälle, Phantasmen genannt — sind, weil sie keine dauernden Ansprüche erheben, von jener Etiquette befreit. — Man achte bei Gesprächen darauf: wie verstimmend, wie unanmuthig werden Meinungssiege erfochten! Wie selten hat Der, der Recht hat, den guten Geschmack, dies verzeihen zu machen!
\par

